%% =====================================================================
%% Replicating a paper --- template
%% ML & FinTech, National Yang Ming Chiao Tung University
%% Instructor: Huei-Wen Teng
%%
%% This file follows the SSRN.tex format used in the lab.
%% Every instruction to you is written in \emph{italics} using \inst{...}.
%% Delete each instruction once you have acted on it.
%% =====================================================================

\documentclass[12pt]{article}

\usepackage{adjustbox,array,amssymb,amsmath,amsfonts,eurosym,geometry,graphicx,
            caption,color,setspace,sectsty,comment,footmisc,natbib,pdflscape,
            subfigure,longtable,lscape,float,multirow,booktabs,makecell,
            hyperref,ulem,appendix}

\normalem
\onehalfspacing

\newtheorem{theorem}{Theorem}
\newtheorem{corollary}[theorem]{Corollary}
\newtheorem{proposition}{Proposition}
\newenvironment{proof}[1][Proof]{\noindent\textbf{#1.} }{\ \rule{0.5em}{0.5em}}

\newtheorem{hyp}{Hypothesis}
\newtheorem{subhyp}{Hypothesis}[hyp]
\renewcommand{\thesubhyp}{\thehyp\alph{subhyp}}

\newcolumntype{L}[1]{>{\raggedright\let\newline\\arraybackslash\hspace{0pt}}m{#1}}
\newcolumntype{C}[1]{>{\centering\let\newline\\arraybackslash\hspace{0pt}}m{#1}}
\newcolumntype{R}[1]{>{\raggedleft\let\newline\\arraybackslash\hspace{0pt}}m{#1}}

\geometry{left=1.0in,right=1.0in,top=1.0in,bottom=1.0in}

%% ---------------------------------------------------------------------
%% Instructions are italic. Delete the \inst{...} once you have done it.
%% ---------------------------------------------------------------------
\newcommand{\inst}[1]{\textit{#1}}

\begin{document}

\begin{titlepage}

\title{Replicating the paper:
       \inst{exact title of the paper you are replicating, followed by the
       citation, e.g.\ \cite{bitetto2024can}}}

\author{\inst{Your name}\thanks{Department of Information Management and Finance,
        National Yang Ming Chiao Tung University. Email:
        \href{mailto:you@nycu.edu.tw}{you@nycu.edu.tw}}}

\date{\today}
\maketitle

\begin{abstract}
\noindent
\inst{Briefly explain what the paper you are replicating does, what it
contributes, and how it bridges the gap left by the papers before it. Then say
why you are interested in this paper. Write it as one paragraph of continuous
prose, not as a list. Aim for 150--200 words. Do not simply copy the original
abstract: the abstract of a replication describes the replication.}

\vspace{0in}
\noindent\textbf{Keywords:} \inst{copy the keywords from the paper you are replicating}

\vspace{0in}
\noindent\textbf{JEL Codes:} \inst{copy the JEL codes from the paper you are
replicating. If the paper does not give any, choose from the JEL classification
system and say in a footnote that you chose them yourself.}

\bigskip
\end{abstract}

\setcounter{page}{0}
\thispagestyle{empty}
\end{titlepage}
\pagebreak

\newpage
\clearpage


\section{Motivation}
\label{sec:motivation}

\inst{Treat this as a short journal article rather than a homework write-up.
The whole manuscript is limited to 2,400 words, excluding tables, figures and
references, so every paragraph has to earn its place. This file already uses the
SSRN format the lab writes in, so if you later develop the work into a thesis or
a conference paper you will not have to move it.}

\inst{Your introduction should do four things, in this order. Say what question
the original paper asks and why it matters. Say what the paper found. Say what
you set out to reproduce, and be explicit about what you could and could not
obtain --- the same data, a subset, or a substitute. Finally, say what a reader
learns from your replication that they would not learn from the paper alone.}

\inst{Cite five papers that are fundamental to the question, and five that
attack it with machine learning or AI. Follow the Google Scholar convention for
bibkeys, as in \texttt{bitetto2024can}. A citation you have not read does not
belong here.}


\section{Data}

\subsection{Dependent and independent variables}

\inst{List every variable with its acronym, a one-line description, and its type: binary, categorical, or continuous. State clearly which variable is the response
and which are the predictors. If your data differs from the data in the original paper, say so here and say what you expect that to change.}


\subsection{Data exploration, cleaning and preprocessing}

\subsubsection{Cross-sectional data}

\begin{itemize}
\item Check whether there are missing values.
  \begin{itemize}
  \item It is not sufficient to list the proportion of missing values per variable.
  \item Show exactly which observations and which variables the missing values fall in.
  \item If values are missing, consider a missing indicator.
  \end{itemize}
\item Check whether there are outliers.
  \begin{itemize}
  \item Check first whether they are errors introduced when the data was retrieved.
        If they are not, keep them for the analysis.
  \end{itemize}
\item For categorical variables, show the frequency tables.
  \begin{itemize}
  \item For predictors, decide whether to use dummy variables or to treat the variable
        as continuous, and whether to merge categories that are extremely rare.
  \item For the response, decide between logistic and multinomial regression, and
        check whether the classes are imbalanced.
  \end{itemize}
\item For continuous variables, report mean, standard deviation, minimum, maximum,
      skewness and kurtosis.
  \begin{itemize}
  \item Decide whether a logarithm is needed, which it usually is when the data is
        heavily skewed.
  \end{itemize}
\item Scaling matters. Min--max or z-score puts the features on a common range and
      makes the numerical procedures more stable. Report whether you scaled or not.
\item Use pairwise scatter plots to look for patterns, and decide whether quadratic
      terms, two-way interactions or logarithms are called for.
\item Show a heatmap of the correlations.
\end{itemize}


\section{Methodology}

\subsection{Benchmark model}

\inst{Use the same benchmark the original paper used, so that your numbers are
comparable to theirs. If you cannot, say which model you substituted and why.}

One of the following:
\begin{itemize}
\item Logistic regression
\item Multinomial regression
\item Linear regression
\end{itemize}


\subsection{Experimental Design}

\inst{Give a flow chart of your procedure, from raw data to reported number.
A reader should be able to reproduce your result from that figure alone.}

\begin{itemize}
\item Explain the design of the replication.
\item State what is novel in your version. For a replication this is usually not a
      new method: it is a test the original paper did not run, data it did not use,
      or an assumption it did not check.
\end{itemize}


\subsection{Performance measures}

\inst{Use the measures the original paper reports, so the comparison is direct.
Add any measure you think the paper should have reported, and say why.}


\section{Results}

\inst{Put your numbers and the original paper's numbers side by side in one table.
Then state plainly whether the result replicated. ``It did not replicate'' is a
finding, not a failure, provided you can show what you did.}


\section{Conclusion}

\inst{Two paragraphs. What survived the replication, and what a reader should now
believe less strongly than before.}


\section*{Declaration}

\subsection*{Declaration of Pedagogical Purpose}

This manuscript is prepared as part of the course project for Machine Learning and
FinTech at National Yang Ming Chiao Tung University, instructed by Professor
Huei-Wen Teng. The work is intended solely for pedagogical purposes, including
learning, discussion, and academic training. It does not represent original research
submitted for journal publication, nor should it be considered professional financial
advice.

\subsection*{Declaration of generative AI and AI-assisted technologies in the writing process}

\inst{Name the tools you actually used. Replace the bracketed placeholders below.}

During the preparation of this work, the authors used [ChatGPT] and [Grok], developed
by [OpenAI] and [xAI] respectively, in order to improve the clarity, precision, and
overall quality of the writing. After using these tools, the authors reviewed and
edited the content as needed and take full responsibility for the content of the
article.


\bibliographystyle{chicago}
\bibliography{ref}


\appendix

\section{\inst{Appendix title}}

\inst{Titles have to be concise. Put here anything a reader needs in order to
reproduce your work but does not need in order to follow the argument.}


\end{document}
